In this work, we presented GEAK, a modular and agent-based framework that leverages inference-time compute scaling with frontier LLMs to automatically generate efficient Triton GPU kernels from text task descriptions. By coordinating multiple specialized agents, GEAK iteratively improves both the correctness and performance of generated code without requiring additional training.

To rigorously assess kernel generation capabilities, we introduced two benchmarks: TritonBench-revised benchmark, a curated subset of TritonBench-G with improved test harnesses, and the newly constructed ROCm benchmark, which consists of real-world kernels from open-source AMD repositories. These benchmarks provide a diverse and challenging testbed for evaluating both functional correctness and runtime efficiency.

By fully open-sourcing the benchmarks and code for running the agent, we expect to engage the open-source community to accelerate the development of GPU kernels. We invite developers, researchers, and AI enthusiasts to explore the agent and benchmarks, and we hope this will foster innovation and collaboration within the AI community to come up with even better methods and final output kernels that can significantly improve the efficiency of training and inference for large-scale AI models.